\documentclass[12pt,a4paper,reqno]{amsart}

% language
\usepackage[greek,english]{babel}
\usepackage[utf8]{inputenc}
\usepackage{alphabeta}

% change default names to greek
\addto\captionsenglish{
    \renewcommand{\contentsname}{Περιεχόμενα}
    \renewcommand{\refname}{Βιβλιογραφία}
    \renewcommand{\datename}{Ημερομηνία:}
    \renewcommand{\urladdrname}{Ιστοσελίδα}
}

% math 
\usepackage{amsmath,amsthm,amssymb,amscd,stmaryrd} 

% font
\usepackage[cal=euler]{mathalfa}
\usepackage{libertinus-type1}
% \usepackage{txfonts} % for upright greek letters
\usepackage{bm} % for bold symbols
\usepackage{bbm} % for the simply-looking bb symbols

% miscellaneous 
\usepackage{changepage} %for indenting environments
\usepackage{csquotes} % example: \textcquote{}
\usepackage{blkarray}
\setcounter{MaxMatrixCols}{20} % default for pmatrix is 10!!
\usepackage{ytableau}
\usepackage{array} %needed to increase the vertical length in a tabular

% drawing
\usepackage{tikz,tikz-cd}
\usetikzlibrary{shapes.misc, patterns, matrix, calc, intersections,positioning,arrows.meta}
\usepackage{graphics,graphicx}
\usepackage{float} % provides enhanced control and customization options for floating objects such as figures and tables

% colors
\usepackage{xcolor}
\definecolor{darkcandyapplered}{rgb}{0.64, 0.0, 0.0}
\definecolor{midnightblue}{rgb}{0.1, 0.1, 0.44}
\definecolor{mylightblue}{HTML}{336699}
\definecolor{burntorange}{rgb}{0.8, 0.33, 0.0}
\definecolor{iceberg}{rgb}{0.44, 0.65, 0.82}
\definecolor{applegreen}{rgb}{0.55, 0.71, 0.0}
\definecolor{canaryyellow}{rgb}{1.0, 0.94, 0.0}
\definecolor{lightgray}{rgb}{0.83, 0.83, 0.83}

% hrefs
\usepackage{hyperref}
\usepackage[noabbrev,capitalize]{cleveref}
\hypersetup{
    pdftoolbar=true,        
    pdfmenubar=true,        
    pdffitwindow=false,     
    pdfstartview={FitH},  % fits the width of the page to the window
    pdftitle={},
    pdfauthor={},
    pdfsubject={},
    pdfkeywords={},
    pdfnewwindow=true,  % links in new window
    colorlinks=true,  % false: boxed links; true: colored links
    linkcolor=darkcandyapplered,   % color of internal links
    citecolor=midnightblue,  % color of links to bibliography
    urlcolor=cyan,  % color of external links
    linktocpage=true  % changes the links from the section body to the page number
    }

% geometry
\textwidth=16cm 
\textheight=21cm 
\hoffset=-55pt 
\footskip=25pt

% thm envs (you might need to change the path)
\theoremstyle{definition}
\newtheorem{theorem}{Θεώρημα}
\newtheorem{proposition}[theorem]{Πρόταση}
\newtheorem{lemma}[theorem]{Λήμμα}
\newtheorem{corollary}[theorem]{Πόρισμα}
\newtheorem{conjecture}[theorem]{Εικασία}
\newtheorem{problem}[theorem]{Πρόβλημα}
\newtheorem*{claim}{Ισχυρισμός}
\newtheorem{observation}[theorem]{Παρατήρηση}
\newtheorem{definition}[theorem]{Ορισμός}
\newtheorem{question}[theorem]{Ερώτημα}
\newtheorem*{questions}{Ερωτήματα}
\newtheorem*{que}{Ερώτημα}
\newtheorem{example}[theorem]{Παράδειγμα}
\newtheorem{exercise}{Άσκηση}
\newtheorem*{zorn}{Λήμμα του Zorn}
\newtheorem*{example_6.6}{Παράδειγμα~6.6}

\theoremstyle{remark}
\newtheorem*{remark}{Παρατήρηση}

% fixes the correct numbering of environments
\numberwithin{theorem}{section}
\numberwithin{exercise}{section}
\numberwithin{equation}{section}

% math ops 
% fields
\renewcommand{\AA}{\mathbbmss{A}} 
\newcommand{\NN}{\mathbbmss{N}} 
\newcommand{\ZZ}{\mathbbmss{Z}} 
\newcommand{\QQ}{\mathbbmss{Q}} 
\newcommand{\RR}{\mathbbmss{R}} 
\newcommand{\CC}{\mathbbmss{C}} 
\newcommand{\KK}{\mathbbmss{K}} 
\newcommand{\FF}{\mathbbmss{F}} 
\newcommand{\HH}{\mathbbmss{H}} 

% frak
\newcommand{\fS}{\mathfrak{S}}  
\newcommand{\fm}{\mathfrak{m}}  
\newcommand{\fp}{\mathfrak{p}}  

% calligraphic 
\newcommand{\aA}{\mathcal{A}} 
\newcommand{\bB}{\mathcal{B}}
\newcommand{\cC}{\mathcal{C}}
\newcommand{\dD}{\mathcal{D}}
\newcommand{\eE}{\mathcal{E}}
\newcommand{\fF}{\mathcal{F}}
\newcommand{\hH}{\mathcal{H}}
\newcommand{\iI}{\mathcal{I}}
\newcommand{\lL}{\mathcal{L}}
\newcommand{\oO}{\mathcal{O}}
\newcommand{\pP}{\mathcal{P}}
\newcommand{\qQ}{\mathcal{Q}}
\newcommand{\sS}{\mathcal{S}}
\newcommand{\mM}{\mathcal{M}}
\newcommand{\uU}{\mathcal{U}}
\newcommand{\vV}{\mathcal{V}}
\newcommand{\zZ}{\mathcal{Z}}

% bold
\newcommand{\bfa}{\pmb{a}}
\newcommand{\bfe}{\mathbf{e}}
\newcommand{\bfF}{\pmb{F}}
\newcommand{\bfR}{\pmb{R}}
\newcommand{\bfv}{\mathbf{v}}
%\newcommand{\bfx}{\bm{x}}
%\newcommand{\bfx}{\mathbf{x}} 
\newcommand{\bfx}{\pmb{x}}
\newcommand{\bfX}{\pmb{X}}
\newcommand{\bfy}{\pmb{y}}
\newcommand{\bfz}{\pmb{z}}

% roman
\newcommand{\rmA}{\mathrm{A}}
\newcommand{\rmB}{\mathrm{B}}
\newcommand{\rmC}{\mathrm{C}}
\newcommand{\rmD}{\mathrm{D}} 
\newcommand{\rmF}{\mathrm{F}} 
\newcommand{\rmH}{\mathrm{H}} 
\newcommand{\rmh}{\mathrm{h}} 
\newcommand{\rmI}{\mathrm{I}} 
\newcommand{\rmK}{\mathrm{K}}
\newcommand{\rmM}{\mathrm{M}}
\newcommand{\rmP}{\mathrm{P}}  
\newcommand{\rmp}{\mathrm{p}}  
\newcommand{\rmQ}{\mathrm{Q}}  
\newcommand{\rmR}{\mathrm{R}}
\newcommand{\rmS}{\mathrm{S}}
\newcommand{\rmT}{\mathrm{T}}
\newcommand{\rmU}{\mathrm{U}}
\newcommand{\rmV}{\mathrm{V}}
\newcommand{\rmY}{\mathrm{Y}}
\newcommand{\rmZ}{\mathrm{Z}}
\newcommand{\rmz}{\mathrm{z}}

% arrows and symbols 
\renewcommand{\to}{\rightarrow}
\newcommand{\toto}{\longrightarrow}
\newcommand{\mapstoto}{\longmapsto}
\newcommand{\then}{\Rightarrow}
\newcommand{\IFF}{\Leftrightarrow}
\newcommand{\tl}{\tilde}
\newcommand{\wtl}{\widetilde}
\newcommand{\ol}{\overline}
\newcommand{\ul}{\underline}
\newcommand{\oldemptyset}{\emptyset}
\renewcommand{\emptyset}{\varnothing}
\DeclareMathSymbol{\Arg}{\mathbin}{AMSa}{"39} % for arguments 
\newcommand{\onto}{\ensuremath{\twoheadrightarrow}}
\newcommand{\tle}{\trianglelefteq}
\newcommand{\tge}{\trianglerighteq}

% custom ops
\newcommand{\rmKer}{\mathrm{Ker}}
\newcommand{\rmIm}{\mathrm{Im}}
\newcommand{\lcm}{\mathrm{lcm}}
\newcommand{\GL}{\mathrm{GL}}
\newcommand{\ev}{\mathrm{ev}}
\newcommand{\End}{\mathrm{End}}
\newcommand{\rmid}{\mathrm{id}}
\newcommand{\rmspan}{\mathrm{span}}
\newcommand{\Hom}{\mathrm{Hom}}
\newcommand{\Ann}{\mathrm{Ann}}
\newcommand{\Set}{\pmb{\mathrm{Set}}}
\newcommand{\Rmod}{\pmb{\mathrm{mod}}}
\newcommand{\rmrank}{\mathrm{rank}}
\newcommand{\mSpec}{\mathrm{mSpec}}
\newcommand{\Spec}{\mathrm{Spec}}
\newcommand{\tor}{\mathrm{tor}}
\newcommand{\rad}{\mathrm{rad}}
\newcommand{\Pol}[1]{F[x_1, x_2, \dots, x_{#1}]}

% absolute value symbol
\usepackage{mathtools} 
\DeclarePairedDelimiter\abs{\lvert}{\rvert}%
\DeclarePairedDelimiter\norm{\lVert}{\rVert}%
\makeatletter
\let\oldabs\abs
\def\abs{\@ifstar{\oldabs}{\oldabs*}}

% tensor symbol
\newcommand{\tensor}[1]{%
  \mathbin{\mathop{\otimes}\limits_{#1}}%
}

% permutation cycle notation
\ExplSyntaxOn
\NewDocumentCommand{\cycle}{ O{\;} m }
 {
  (
  \alec_cycle:nn { #1 } { #2 }
  )
 }

\seq_new:N \l_alec_cycle_seq
\cs_new_protected:Npn \alec_cycle:nn #1 #2
 {
  \seq_set_split:Nnn \l_alec_cycle_seq { , } { #2 }
  \seq_use:Nn \l_alec_cycle_seq { #1 }
 }
\ExplSyntaxOff

% setminus symbol
\newcommand{\mysetminusD}{\hbox{\tikz{\draw[line width=0.6pt,line cap=round] (3pt,0) -- (0,6pt);}}}
\newcommand{\mysetminusT}{\mysetminusD}
\newcommand{\mysetminusS}{\hbox{\tikz{\draw[line width=0.45pt,line cap=round] (2pt,0) -- (0,4pt);}}}
\newcommand{\mysetminusSS}{\hbox{\tikz{\draw[line width=0.4pt,line cap=round] (1.5pt,0) -- (0,3pt);}}}
\newcommand{\sm}{\mathbin{\mathchoice{\mysetminusD}{\mysetminusT}{\mysetminusS}{\mysetminusSS}}}

% miscellaneous commands
\newcommand{\defn}[1]{{\color{mylightblue}{#1}}}
\newcommand{\toDo}{{\bf\color{red} TODO}}
\newcommand{\toCite}{{\bf\color{green} CITE}}
\newcommand*{\vertbar}{\rule[-1ex]{0.5pt}{2.5ex}} % for matrices with column vectors
\newcommand*{\horzbar}{\rule[.5ex]{2.5ex}{0.5pt}} % for matrices with row vectors
\newcommand{\myblue}[1]{{\color{iceberg}{#1}}}
\newcommand{\myorange}[1]{{\color{burntorange}{#1}}}
\newcommand{\mygreen}[1]{{\color{applegreen}{#1}}}
\newcommand{\myred}[1]{{\color{darkcandyapplered}{#1}}}

% ferrer's diagram
\newcommand{\fdiagram}[1]{
    \begin{tikzpicture}[scale=.7]
        \fill foreach \Z [count=\Y] in {#1}
        {foreach \X in {1,...,\Z} 
        {(\X,-\Y) circle[radius=3pt]}};
    \end{tikzpicture}
}

%
\usepackage{pgfplots}
\usepgfplotslibrary{groupplots}
\pgfplotsset{compat=1.18}

%
\makeatletter
\def\mathcolor#1#{\@mathcolor{#1}}
\def\@mathcolor#1#2#3{%
  \protect\leavevmode
  \begingroup
    \color#1{#2}#3%
  \endgroup
}
\makeatother

% restriction
\newcommand\restr[2]{{% we make the whole thing an ordinary symbol
  \left.\kern-\nulldelimiterspace % automatically resize the bar with \right
  #1 % the function
  \littletaller % pretend it's a little taller at normal size
  \right|_{#2} % this is the delimiter
  }}
\newcommand{\littletaller}{\mathchoice{\vphantom{\big|}}{}{}{}}

% 
\newenvironment{nouppercase}{%
  \let\uppercase\relax%
  \renewcommand{\uppercasenonmath}[1]{}}{}

% titlepage
\title{226: Θεωρία Δακτυλίων και Modules}
\author[Β.~Δ. Μουστακας]{Βασίλης Διονύσης Μουστάκας \\ Πανεπιστήμιο Κρήτης}
\date{19 Μαΐου 2026}
% \urladdr{\href{https://sites.google.com/view/vasmous}{https://sites.google.com/view/vasmous}}

\begin{document}

\begingroup
\def\uppercasenonmath#1{} % this disables uppercase title
\let\MakeUppercase\relax % this disables uppercase authors
\maketitle
\endgroup

\setcounter{section}{11}
\setcounter{theorem}{4}
\setcounter{equation}{2}
\begin{center}
    \textbf{11. Το Θεώρημα μηδενικών του Hilbert
} (Συνέχεια)
\end{center}

Θα εμπλουτίσουμε την \textquote{γέφυρα} μεταξύ άλγεβρας και γεωμετρίας που \textquote{έχτισε} το ισχυρό θεώρημα μηδενικών του Hilbert συζητώντας κάποιες πολλαπλότητες που μοιάζουν να είναι τα δομικά συστατικά όλων των πολλαπλοτήτων.
  
Η Πρόταση 9.3 μας πληροφόρησε ότι η ένωση δύο πολλαπλοτήτων είναι πολλαπλότητα. Για παράδειγμα, 
\[
\vV(\{xz, yz\}) = \vV(\{x, y\}) \cup \vV(z) \ \subseteq \AA_F^3
\]
Με άλλα λόγια, η πολλαπλότητα αυτή αποτελείται από δυο κομμάτια, μια \textquote{ευθεία} και ένα \textquote{επίπεδο}, τα οποία διαισθητικά μοιάζουν να είναι \textquote{αδιάσπαστα}.

\begin{definition}
  \label{def:variety_irreducible}
  Μια (μη κενή) πολλαπλότητα $V \subseteq \AA_F^n$ ονομάζεται\footnote{Αυτός είναι ο ορισμός της πολλαπλότητας στα συγγράμματα όπου χρησιμοποιείται ο όρος αλγεβρικός σύνολο γι' αυτό που ονομάσαμε εμείς πολλαπλότητα.} \defn{ανάγωγη} (irreducible) αν δεν μπορεί να γραφεί ως $V = V_1 \cup V_2$, για κάποιες γνήσιες πολλαπλότητες $V_1, V_2 \subsetneqq V$.
\end{definition}

Για παράδειγμα, η πολλαπλότητα $\vV(\{xz, yz\})$ δεν είναι ανάγωγη. Διαισθητικά, ένα ση\-μείο, μια ευθεία και ένα επίπεδο θα πρέπει να είναι ανάγωγα. Με γνώμονα την διαίσθηση, το ίδιο πρέπει να ισχύει και για την πολλαπλότητα $\vV(y-x^2)$ στο πραγματικό αφφινικό χώρο. Αλλά, πως μπορούμε να το αποδείξουμε;

\begin{proposition}
  \label{prop:prime_irreducible}
  Μια πολλαπλότητα είναι ανάγωγη αν και μόνο αν το ιδεώδες μηδενισμού της είναι πρώτο, ή ισοδύναμα, αν και μόνο αν ο δακτύλιος συντεταγμένων της είναι ακέραια περιοχή.
\end{proposition}

\begin{proof}[Απόδειξη]
  Ο δεύτερος ισχυρισμός έπεται άμεσα από την Πρόταση 1.11. Για τον πρώτο ισχυρι\-σμό, έστω $V$ μια ανάγωγη πολλαπλότητα. Θα δείξουμε ότι το $\iI(V)$ είναι πρώτο ιδεώδες. Αν $fg \in \iI(V)$, τότε 
  \[
  V = \left(
    V \cap \vV(f)
  \right) \cup \left(
    V \cap \vV(g)
  \right),
  \]
  διότι για κάθε $\bfa \in V$ έχουμε $f(\bfa)g(\bfa) = 0$ που σημαίνει ότι είτε $f(\bfa) = 0$, είτε $g(\bfa) = 0$. Όμως, επειδή η $V$ είναι ανάγωγη, έπεται ότι είτε $V = V \cap \vV(f)$, είτε $V = V \cap \vV(g)$. Κατά συνέπεια, είτε $f \in \iI(V)$, είτε $g \in \iI(V)$.
  
  Αντιστρόφως, υποθέτουμε ότι το $\iI(V)$ είναι πρώτο ιδεώδες και ότι $V = V_1 \cup V_2$, για κάποιες πολλαπλότητες $V_1, V_2 \subseteq V$. Αν $V \neq V_1$, τότε ισχυριζόμαστε ότι $\iI(V) = \iI(V_2)$ και κατά συνέπεια $V = V_2$, διότι η απεικόνιση $W \mapsto \iI(W)$ είναι ένα-προς-ένα (γιατί;). Πράγματι, επειδή $V_2 \subseteq V$, έπεται ότι $\iI(V) \subseteq \iI(V_2)$ (βλ. παρατήρηση μετά τον Ορισμό 9.6). Για τον αντίστροφο εγκλεισμό, έστω $f \in \iI(V_2)$. Επειδή $V_1 \subsetneqq V$ από την ίδια παρατήρηση έπεται ότι $\iI(V) \subsetneqq \iI(V_1)$ και γι αυτό μπορούμε να διαλέξουμε $g \in \iI(V_1) \sm \iI(V)$. Όμως, $V = V_1 \cup V_2$ και γι αυτό $fg \in \iI(V)$ που σημαίνει ότι είτε $f \in \iI(V)$, είτε $g \in \iI(V)$, διότι το $\iI(V)$ είναι πρώτο ιδεώδες. Το δεύτερο δεν μπορεί να συμβεί εξ' επιλογής του $g$ και γι αυτό $f \in \iI(V)$ και το ζητούμενο έπεται.
\end{proof}

Εφαρμόζοντας την Πρόταση \ref{prop:prime_irreducible}, μπορούμε να δείξουμε ότι οι πολλαπλότητες 
\[     
\vV(\{x,y\}), \ \vV(z), \ \text{και} \ \vV(y-x^2)
\]
είναι ανάγωγες. Για την πρώτες δύο, παρατηρούμε 
\[
\frac{F[x,y,z]}{(x,y)} \cong F[z], \quad \text{και} \quad 
\frac{F[x,y,z]}{(z)} \cong F[x,y],
\]
αντίστοιχα, (γιατί;) τα οποία είναι ακέραιες περιοχές. Για την τρίτη, από το Θεώρημα 10.3 έπεται ότι 
\[
\iI\left(\vV\left(y-x^2\right)\right) = \sqrt{(y-x^2)} = (y-x^2).
\]
Όμως, 
\[
\frac{F[x,y]}{(y-x^2)} \cong F[x],
\]
το οποίο είναι επίσης ακέραια περιοχή.

Συνοψίζουμε τα αποτελέσματα της ενότητας αυτής στο παρακάτω \textquote{αλγεβρογεωμετρικό λεξικό}, με την προϋπόθεση ότι δουλεύουμε πάνω από αλγεβρικά κλειστό σώμα:
\begin{table}[H]
\centering
\renewcommand{\arraystretch}{1.3} 
\begin{tabular}{c|c}
\textbf{Άλγεβρα}       & \textbf{Γεωμετρία} \\ \hline
ριζικό ιδεώδες         & πολλαπλότητα \\ \hline 
πρώτο ιδεώδες          & ανάγωγη πολλαπλότητα \\ \hline 
μεγιστικό ιδεώδες      & σημεία του αφφινικού χώρου \\ \hline 
πρόσθεση ιδεωδών       & τομή πολλαπλοτήτων \\ \hline 
γινόμενο/τομή ιδεωδών  & ένωση πολλαπλοτήτων 
\end{tabular}
\end{table}
\noindent

\newpage
\setcounter{section}{12}
\setcounter{theorem}{0}
\setcounter{equation}{0}
\begin{center}
    \textbf{12. Το Λήμμα κανονονικοποίησης της Noether}
\end{center}

TBW.


% \begin{thebibliography}{99}
%     \bibitem{Alu09}
%     P.~Aluffi,
%     \textit{Algebra, Chapter 0},
%     Grad. Stud. in Math. {\bf104},
%     Amer. Math. Soc., Providence, RI, 2009.
%     \bibitem{DF04}
%     D. S. Dummit και R. M. Foote, 
%     \textit{Abstract algebra},
%     third edition, Wiley, 2004. 
% \end{thebibliography}
\end{document}